% Assignment source: ne630/homework/markdown/HW12.md.
% Legacy foundation: ne630_problems/lesson_12/statement.tex and hw12_sol.ipynb.

\noindent{\bf Problem 1}. Equations (3.50)--(3.53) in FNRP describe how to
compute the resonance integral. Use these equations with nuclear data from
OpenMC to compute the absorption resonance integral, $I_a$, for
\isoA{U}{238} over $1\le E\le10^5~\mathrm{eV}$.
\begin{enumerate}[label=(\alph*)]
\item Compute $I_a$.
\item How does this value compare with the value in FNRP Table 3.2?
\item What is the corresponding effective microscopic absorption cross
      section, $\bar{\sigma}_a$, for a $1/E$ flux spectrum over the same
      energy interval?
\end{enumerate}
Use absorption, including radiative capture and fission, rather than
radiative capture alone. Report the resonance integral and effective cross
section in barns.

\begin{adaptedsolution}
The resonance integral and corresponding effective cross section are
\[
 I_a=\int_1^{10^5}\frac{\sigma_a(E)}{E}\,dE,
 \qquad
 \bar{\sigma}_a=\frac{I_a}{\int_1^{10^5}E^{-1}\,dE}
 =\frac{I_a}{\ln(10^5)}.
\]
Using OpenMC pointwise data with
$\sigma_a=\sigma_\gamma+\sigma_f$:
\begin{minted}{python}
import openmc
import numpy as np

u238 = openmc.data.IncidentNeutron.from_hdf5(data_path + "U238.h5")
sig_g = u238.reactions[102].xs["294K"]
sig_f = u238.reactions[18].xs["294K"]

E = np.logspace(0, 5, 100000)
I_a = np.trapz((sig_g(E) + sig_f(E))/E, E)
sig_a_bar = I_a / np.log(1e5)
print(I_a, sig_a_bar)
\end{minted}
The legacy solution notebook gives
\[
 \boxed{I_a\approx2.75\times10^2~\mathrm b},
 \qquad
 \boxed{\bar{\sigma}_a\approx23.9~\mathrm b}.
\]
The integral is close to the FNRP Table~3.2 value of about
$278~\mathrm b$.
\end{adaptedsolution}

\newpage

\noindent{\bf Problem 2}. A major approximation in Eqs.~(3.50)--(3.53) is the
explicit assumption that $\phi(E)=1/E$. Instead, use the same nuclear data as
in Problem 1, supplemented by the total cross sections needed for the mixture,
with the narrow-resonance (NR) spectrum
\[
 \phi_{\mathrm{NR}}(E)\propto\frac{1}{\Sigma_t(E)E}.
\]
Consider an infinite, homogeneous 100-to-1 atom mixture of \isoA{H}{1} and
\isoA{U}{238}, where $\Sigma_t(E)$ is the total macroscopic cross section of
the mixture.
\begin{enumerate}[label=(\alph*)]
\item Compute the effective microscopic absorption cross section,
      $\bar{\sigma}_a$, for \isoA{U}{238} over $1\le E\le10^5~\mathrm{eV}$,
      using the NR spectrum as the flux weighting.
\item Explain why this value differs from the value computed in Problem 1.
\end{enumerate}

\begin{adaptedsolution}
The flux-weighted effective cross section is
\[
 \bar{\sigma}_a
 =\frac{\int_1^{10^5}\sigma_a^{238}(E)\phi_{\mathrm{NR}}(E)\,dE}
        {\int_1^{10^5}\phi_{\mathrm{NR}}(E)\,dE}.
\]
Because the \isoA{U}{238} number density cancels, the mixture total can be
formed from microscopic cross sections as
\[
 \Sigma_t(E)\propto
 100\sigma_t^{\isoA{H}{1}}(E)+\sigma_t^{\isoA{U}{238}}(E).
\]
The OpenMC-data calculation is
\begin{minted}{python}
h1 = openmc.data.IncidentNeutron.from_hdf5(data_path + "H1.h5")
sig_s_u238 = u238.reactions[2].xs["294K"]
sig_g_h1 = h1.reactions[102].xs["294K"]
sig_s_h1 = h1.reactions[2].xs["294K"]

sig_t_u238 = sig_g(E) + sig_f(E) + sig_s_u238(E)
sig_t_h1 = sig_g_h1(E) + sig_s_h1(E)
phi_nr = 1.0 / ((100*sig_t_h1 + sig_t_u238)*E)
sig_a_bar = np.trapz((sig_g(E) + sig_f(E))*phi_nr, E) / np.trapz(phi_nr, E)
\end{minted}
The legacy solution notebook gives
\[
 \boxed{\bar{\sigma}_a\approx8.53~\mathrm b}.
\]
This is much smaller than the $1/E$ result because the NR spectrum includes
energy self-shielding.  The flux is depressed near \isoA{U}{238} absorption
resonances, so those large cross sections receive less flux weight than they
do under an unshielded $1/E$ spectrum.
\end{adaptedsolution}

\clearpage
